\chapter{Related Work}\label{chap:related-work}

This chapter reviews advances in congestion control.
Section~\ref{sec:rw-tcp} covers the improvements proposed for congestion control in
TCP. Section~\ref{sec:rw-quic} covers the
congestion control work done for QUIC, and separates algorithms that are only ported over from TCP from those that use features TCP does not have.

\section{Advances in congestion control}\label{sec:rw-tcp}

\subsection{Improving startup}\label{sec:rw-startup}

Slow start begins with the initial window and doubles the congestion window every RTT,
so it needs about $\log_2(\mathit{BDP}/\mathit{IW})$ rounds to fill a path. This causes two problems. A flow that ends within these
rounds never uses the capacity of the path, and a flow that continues overshoots it:
the window keeps doubling until the first loss is detected, about one RTT after it
occurred, so many packets can be lost at once~\cite{ha2011taming}.

Flach et al.\ analysed billions of
TCP connections to Google's services, and found that connections without
loss complete close to the ideal latency of one RTT, while transfers with loss take
five times longer on average, and nearly 10\% of the connections saw at least one
loss~\cite{flach2013reducing}. Loss is expensive mainly because of how it is
recovered: 77\% of the losses were recovered only by a retransmission timeout rather
than by fast recovery, and the median timeout was six times the
RTT.


\paragraph{Better exit.} A first group of algorithms keeps the exponential growth of
slow start and changes only when it ends. Leaving slow start before the window
overshoots the BDP avoids a burst of losses, but the window grows no faster than
before.

Hybrid Slow Start (HyStart)~\cite{ha2008hystart,ha2011taming} uses two signals that indicate when the window is
approaching the BDP. The first is the length of the
ACK train of a round. The ACKs for a window sent back to back return spaced by
the bottleneck, so the time between the first and the last ACK of the train grows
with the window. The second is the RTT: if the RTTs measured at the start of a round
rise above the minimum RTT seen so far, a queue is building. When either signal fires,
HyStart sets \textit{ssthresh} to the current window and moves to congestion
avoidance, without waiting for a loss. Linux has enabled HyStart by default with CUBIC
since kernel 2.6.29. In practice, the ACK train signal proved unreliable: ACK compression distorts the spacing of the ACKs, so it often fails to detect congestion early~\cite{rfc9406}.

HyStart++~\cite{rfc9406} keeps only the delay
signal. When the minimum RTT of a round exceeds that of the previous
round by more than an eighth, with the threshold clamped between 4 and 16\,ms, the
sender enters Conservative Slow Start, in which the window grows at a quarter of the
rate of slow start. If the RTT falls back, the rise is treated as spurious and slow
start resumes; otherwise, the sender moves to congestion avoidance after five rounds.

SEARCH~\cite{kachooei2023search} replaces the delay signal with the delivery rate. It
tracks the delivery rate over a sliding window and exits slow start when the rate has
not increased by the amount that a doubling window should produce. It was designed for
satellite links, where slow start often exits too early, which limits throughput, or
too late, which causes losses. 

\paragraph{Faster growth.} A second group keeps the structure of slow start, but lets
the window grow faster than doubling while the flow is still far from the capacity of
the path.

Paced Chirping~\cite{misund2019paced} sends \emph{chirps}, trains of packets
with decreasing gaps between them, to estimate the available capacity
quickly. The sender lets the queue drain between chirps, so the
queue never holds more than a few packets, which avoids the overshoot of slow start.
Because the gaps shrink within a chirp, its packets probe increasing rates, and the
sender estimates the available capacity from where their queueing delay starts to
rise. It spreads several chirps over each RTT and can probe a
rate without committing to send at it, so it builds confidence in its estimate before
it raises its window, and it leaves start-up once the data in flight approaches the
estimated BDP.

SUSS~\cite{arghavani2024suss} is a sender-side addition to CUBIC's slow
start. At the end of each round, it uses the ACK train it has just received to
predict whether HyStart's conditions for continuing exponential growth will still
hold in the next round. If they will, it sends additional data within the current
round, which raises the growth factor of the window to four, typically during only
the first one or two rounds. The extra packets are paced, so that they do not create
bursts or disturb the ACK-train measurements that HyStart relies on, and the
acceleration is used only while the window is still much smaller than the flow's
fair share. In real-world experiments, SUSS reduced the FCT by more than 20\% for
flows smaller than 5\,MB on paths with RTTs above 50\,ms.

% An independent replication on AWS found average improvements of 27--44\% on wired paths, also with concurrent background traffic, but showed that the RTT alone
% does not predict how much SUSS helps. Applied to BBR's startup, which is already more
% aggressive, the same principle gave only 1--9\% improvement~\cite{gatti2026replicating}.

Nimble Start~\cite{nimblestart} ramps up using the advertised receive window, but never
beyond its current estimate of the BDP. It splits each round into two halves: a slower
first half and a faster second half, which measures the new estimated BDP from the returning ACK train. It leaves startup
when the estimate stops growing notably for a few rounds. To distinguish random losses from congestion
losses, it uses BBR's criterion, a loss rate above 2\% with at least six discontiguous
losses in a round. After a congestion loss, it scales its estimate by 0.7 and keeps
probing, matching the Cubic and BBR configurations. Nimble Start manages to complete short and mid-length flows faster than Cubic, BBR, SUSS and Paced Chirping, benefiting most with large BDPs, while causing no packet loss when the flow is alone.

\paragraph{Receiver-informed probing.} One design moves part of the
measurement to the receiver. In Receiver-Optimized Startup (ROS)~\cite{ros}, the sender
marks probe packets with a header bit, and the receiver measures their arrival rate
over probe windows of 20\,ms. Measured at
the receiver, the rate is not distorted by suppressed or delayed acknowledgments. The
receiver returns its estimate through the advertised receive window field. After
the initial window, the sender enters a cautious phase at 40\% of this first estimate
and keeps probing until the increase stagnates or a congestion loss occurs.

\paragraph{Bypassing slow start.} Another group of algorithms skips slow start
altogether and sends at a high rate from the first RTT. Jump
Start~\cite{liu2007jumpstart} paces all the data it has queued, up to the receiver's
advertised window, over the RTT measured during the handshake, and leaves the rest of
TCP's congestion control unchanged. Its authors note that this
mainly helps medium-sized transfers, which take longer than one RTT but are short
enough that slow start makes up a large part of their lifetime. The price is more
loss: in their simulations, Jump Start's loss rate was more than four times that of
slow start without competing connections, and about twice as high when connections
competed.

Halfback~\cite{li2015halfback} takes the same approach, but tries to mitigate the packet loss caused by the aggressive start. It also begins with an aggressive, paced first round, but in
the following RTT, it proactively retransmits packets to protect itself against
losses, which lets it recover from them quickly. Its authors
report lower latency than any of the earlier schemes they compared, while harming
competing flows significantly less, and Halfback is TCP-friendly and needs changes
only at the sender. The cost is overhead: it retransmits nearly
half of the packets of its startup phase.



\paragraph{Reusing state from earlier connections.} A connection can also skip part of
startup by reusing what earlier connections have learned. RFC~9040~\cite{rfc9040} describes how a TCP
host can share connection state. With \emph{temporal sharing}, a new
connection to a host initializes its RTT, RTT variance, slow start threshold and
congestion window from a cache of earlier connections. With \emph{ensemble sharing},
concurrent connections share this state.

Stateful-TCP~\cite{guo2020stateful} goes further and reuses the bandwidth itself: it
uses the bandwidth that a previous flow to the same destination measured to ramp up
the next flow immediately, without probing in slow start. Applied to CUBIC, it reduced
the FCT by 37.5\% and increased throughput by more than 50\% in Internet experiments
with more than 1\,000 clients.


\subsection{Model-based congestion control}\label{sec:rw-model}

\paragraph{Limits of loss-based control.} Loss-based algorithms such as Reno and CUBIC
reduce their rate only when packets are lost, which has
two consequences. When bottleneck buffers are large, the
algorithms keep them full and add queueing delay to every packet that shares the
bottleneck, a problem known as bufferbloat~\cite{gettys2012bufferbloat}. When buffers
are small, or losses are not caused by congestion, they reduce their rate
unnecessarily. The BBR specification provides an example for the second case: for CUBIC to sustain
10\,Gbit/s over a 100\,ms RTT, the loss rate must stay below 0.000003\%, and a loss
rate of 1\% on the same RTT is at most 3\,Mbit/s~\cite{cardwell2026bbrdraft}.

\paragraph{BBR.} BBR~\cite{cardwell2016bbr} (Bottleneck Bandwidth and Round-trip propagation time) builds an
explicit model of the path from two estimates: the maximum recent delivery rate, which
estimates the bottleneck bandwidth, and the minimum recent RTT, which estimates the
propagation delay. It paces its packets at a rate derived from
the bandwidth estimate and limits the data in flight to a multiple of the estimated
BDP, so that it fills the path without building a long queue. The current IETF draft
specifies the third version, BBRv3, and is
intended for publication as an Experimental RFC~\cite{cardwell2026bbrdraft}.

BBR's own startup still has exponential growth. In its Startup state, the sender uses
a pacing gain of $4\ln 2 \approx 2.77$, the smallest gain that doubles the sending rate
every round, and needs $O(\log_2 \mathit{BDP})$ rounds to find the
bandwidth. It leaves Startup when the delivery rate has
grown by less than 25\% for three rounds, or when the loss rate of a round exceeds 2\%
with at least six loss events~\cite{cardwell2026bbrdraft}.

BBR is widely deployed: it is the default congestion control of YouTube, and Google
has used it for a large share of its traffic, both in its data centers and on its WAN,
since 2017~\cite{ware2019modeling,zeynali2024promises}. Its original publication
showed that it can outperform CUBIC in some environments~\cite{cardwell2016bbr}, and
independent studies have since tested it more. Cao et al.\ found that BBR achieves
higher throughput than loss-based algorithms on paths with shallow buffers, despite
retransmitting much more, while loss-based algorithms such as CUBIC do better with deep
buffers~\cite{cao2019when}. Fairness towards CUBIC has remained
its main weakness. Measurements showed a single BBR flow taking a fixed 40\% of the
capacity even when competing with up to 16 CUBIC or Reno flows~\cite{ware2019modeling}. Google has revised BBR twice in response, but even BBRv3
struggles to share bandwidth evenly with CUBIC across a wide range of network
conditions, and all versions of BBR share bandwidth unfairly between flows whose RTTs
differ substantially~\cite{zeynali2024promises}.

\paragraph{Learning-based congestion control.} Several designs replace classic congestion
control rules with learned or optimized ones. Remy~\cite{winstein2013remy} generates the rules offline: the
designer specifies prior assumptions about the network and an objective, such as high
throughput and low queueing delay, and Remy produces endpoint rules that try to
achieve it. In ns-2 simulations~\cite{ns2}, the generated algorithms
outperformed CUBIC, Compound~\cite{tan2006compound} and Vegas~\cite{brakmo1994vegas}, but their performance depended on how well the
real conditions matched the assumptions given at design time.

PCC~\cite{dong2015pcc} learns online
instead: the sender continuously runs short experiments at different rates and moves
towards the rate that yields a higher value of a utility function of throughput and
loss. PCC Vivace~\cite{dong2018vivace} builds upon PCC by adding latency to the utility function. Pantheon~\cite{yan2018pantheon} provides a testbed of emulated paths
calibrated against real Internet paths, and was used to train Indigo, a controller that
learns from examples.

Aurora~\cite{jay2019aurora} applies deep reinforcement learning
to the PCC framework. In each monitoring interval, a small neural network maps a
history of statistics from recent intervals: the gradient of the latency, the
ratio of the current latency to the minimum, and the ratio of sent to acknowledged
packets, to a change of the sending rate. It is trained in simulation with a reward
that favors throughput and penalizes latency and loss. Although Aurora was trained
only on links of 1.2 to 6\,Mbit/s, it nearly fully used links from 1 to
128\,Mbit/s, with latency far below CUBIC and similar to PCC Vivace, and it kept close
to full throughput at higher random loss rates than both.

Orca~\cite{abbasloo2020orca} combines CUBIC
with a reinforcement-learning agent that adjusts the window at a coarser timescale,
to reduce overhead and to perform well in network conditions not seen during
training. Sage~\cite{yen2023sage} learns a policy offline from the behavior of
existing heuristic algorithms, and Astraea~\cite{liao2024astraea} uses multi-agent
reinforcement learning to optimize fairness, convergence speed and stability between
competing flows. These designs optimize the control of the
sending rate over the lifetime of a flow, and many of them decide once per monitoring
interval, which lasts on the order of an RTT. 

Learning-based congestion control also has drawbacks that have so far limited its use.
A learned policy optimizes its own objective,
which can disrupt fairness, fast convergence and stability between competing
flows~\cite{liao2024astraea}. The authors of Aurora warn that their agent might not be
friendly to TCP, and that an agent trained against TCP could even learn to cause
losses on purpose, to make TCP back off~\cite{jay2019aurora}. Additionally, the decisions of
a neural network are hard to interpret, which makes it difficult to show that a
learned controller is safe to deploy. Formal verification of Aurora's policy~\cite{eliyahu2021verifying} has
exposed concrete scenarios in which basic requirements are violated. Under excellent
network conditions, with latency close to its minimum and no loss, the policy could
keep decreasing its sending rate until it reached the minimum and stayed there. On a
link with a shallow buffer and high, fluctuating loss, it would keep its sending rate
unchanged instead of reducing it.


\subsection{Network-assisted congestion control}\label{sec:rw-network}

Network-assisted algorithms use signals from the routers or switches on the path,
instead of or in addition to loss and delay.

\paragraph{Congestion control for data centres} DCTCP~\cite{rfc8257} reduces its window in
proportion to the fraction of packets that switches mark, instead of reacting to a mark
as to a loss. This keeps queues short in data centers, but DCTCP must not be used on
the public Internet without additional measures, because it takes a larger share of
the capacity than conventional TCP.

HPCC~\cite{li2019hpcc} sets the
window from the queue length and link load that switches add to each packet, so that
flows can start at line rate while queues stay near zero. Bolt~\cite{arslan2023bolt}
uses programmable switches to react to congestion in less than an RTT and to hand the
bandwidth of finishing flows to others at once, which reduced the 99th-percentile FCT
by up to three times compared with HPCC.

\paragraph{Richer explicit feedback.} ABC~\cite{goyal2020abc} lets routers on wireless
links mark each packet as accelerate or brake within the existing ECN field, to
steer senders to a target rate. On Wi-Fi, it achieved 30--40\% higher throughput than
CUBIC with CoDel at similar delays.

\paragraph{Router-assisted startup.} Quick-Start~\cite{rfc4782} lets the routers on a path approve a
fast start directly. The sender puts a requested rate in an IP option,
and each router on the path approves, reduces or rejects it. If all routers approve,
the sender may send at up to the approved rate for one window of data, after which its
normal congestion control resumes. A request that crosses a router, IP tunnel or MPLS
path that does not support Quick-Start is detected, and the sender then uses its
default congestion control.

\paragraph{Limits for Internet paths.} All of these designs depend on support from the
network. DCTCP, HPCC and Bolt run only in data centers, ABC needs ABC routers at the
wireless bottleneck to give its benefits, and Quick-Start needs every router on the path to be compliant.


% \subsection{Coordinating flows and low-priority congestion control}\label{sec:rw-coordination}

% \paragraph{Sharing congestion state between connections.} The Congestion Manager~\cite{balakrishnan1999cm} is a
% module in the host that manages congestion for all of its applications and
% transports. It combines a window-based controller, a
% scheduler that divides the bandwidth among flows, and feedback from receivers, and
% applications use it through an API. Connections to the same destination share
% congestion state, which lets a new connection use what earlier connections have
% learned instead of going through slow start again. ctrlTCP\cite{islam2018ctrltcp}
% couples the congestion controllers of the TCP connections leaving one sender, also
% when they go to different destinations through a shared bottleneck, such as a home
% uplink~. It lets applications set precise shares of the
% bandwidth, and it reduces queueing delay, loss and the completion time of short flows.
% MPAT~\cite{singh2004mpat} keeps one ordinary AIMD control loop per flow, so that a group of $N$ flows takes
% the fair share of $N$ TCP flows, and then shifts acknowledgments and bandwidth between
% their windows to give each flow a chosen share.

% \paragraph{Scavengers.} A scavenger is a congestion control for background traffic,
% which only uses capacity that other flows leave unused. TCP-LP~\cite{kuzmanovic2003tcplp} uses one-way delay as
% an early congestion signal, before TCP would see a loss; the sender then halves its
% window and backs off further if the signal repeats. TCP Nice~\cite{venkataramani2002nice}
% extends TCP Vegas with a more sensitive congestion detector, a multiplicative
% reduction when the RTT rises, and windows smaller than one packet, and it provably
% bounds the interference with foreground flows in a simplified network
% model. LEDBAT~\cite{rfc6817} measures one-way delay and keeps the
% queueing delay it causes below a target of at most 100\,ms, so it yields to standard
% TCP flows. PCC Proteus~\cite{meng2020proteus} observes that LEDBAT often fails to yield and uses
% the deviation of latency as a sign of competition; within the PCC framework, it can act
% as a scavenger, as a primary protocol, or in hybrid modes between the
% two.


\section{Congestion control in QUIC}\label{sec:rw-quic}

\subsection{TCP algorithms ported to QUIC}\label{sec:rw-quic-ports}

RFC~9002~\cite{rfc9002} specifies that
sender may use any congestion controller algorithm, provided that it follows the guidelines for
congestion control over UDP~\cite{rfc8085}. In practice, most QUIC stacks
re-implement standard TCP algorithms, mainly CUBIC, Reno and BBR~\cite{mishra2023cambrian}. A port from TCP is not exact however: QUIC's acknowledgments, loss detection and timers differ
from TCP's, so each stack has to decide how to map
the rules of an algorithm onto them, and each also chooses its own parameters, pacing
and defaults. Measurements show that these choices matter:
\begin{itemize}
  \item Kakhki et al.~\cite{kakhki2017long} found that Google's early QUIC, using CUBIC, took about twice the
    bandwidth of a competing TCP CUBIC flow, because its CUBIC implementation grew the
    window more aggressively and more often.
  \item Marx et al.\cite{marx2020same} compared 15 IETF QUIC implementations and found large differences
    in flow control, congestion control, 0-RTT, multiplexing and packetization, and
    concluded that most implementations were not yet fully optimized or
    validated~.
  \item Mishra et al.~\cite{mishra2023cambrian, mishra2022speciation} developed the first congestion control benchmarking tool for QUIC stacks, QUICbench and compared the throughput and delay of QUIC implementations of
    standard algorithms with the Linux kernel implementations and found significant
    deviations. In 11 stacks, most
    implementations of CUBIC, Reno and BBR behaved like the kernel versions in shallow
    buffers but less so in deep buffers. Five implementations deviated strongly, which
    can make them unfair to other flows, and some deviations came from TCP features
    outside the algorithm itself, such as HyStart, that the QUIC stacks did not
    implement.
  \item Kempf et al.~\cite{kempf2025quicsteps} showed that QUIC stacks implement pacing differently, and that
    user-space timers, system call overhead and operating system scheduling make
    precise pacing difficult.
  \item Mücke et al.~\cite{mucke2025waiting} observed that large providers configure significantly different
    maximum numbers of retransmissions and timeouts in their QUIC
    deployments.
  \item Sidhu and Bentaleb~\cite{sidhu2026video} found that the same congestion control algorithm in
    different QUIC implementations leads to significantly different quality of
    experience for video streaming.
\end{itemize}
The behavior and performance of a congestion control mechanism in QUIC therefore depends on the stack that it runs in, and a
comparison of algorithms is only meaningful within the same stack.

\paragraph{TCP ideas implemented in QUIC.}  Because QUIC runs in user space, new congestion control can be
implemented and tested in a QUIC stack far more easily than in an operating system
kernel. Most of these algorithms, however, are
designed around signals that TCP provides as well.

Two TCP startup mechanisms have been brought to QUIC.
SEARCH~\cite{cronin2024search} was implemented in an open-source QUIC library and evaluated over a geostationary satellite link, where it exited
slow start before any loss in most cases and improved throughput compared with the
baseline. Fischer et al.\ implemented Careful Resume~\cite{fischer2025neqo} in Mozilla's
Neqo stack and evaluated it on emulated paths and on real fiber, 5G and satellite
paths. It reduced the time to the last byte by several RTTs, and by 12 RTTs on a
1\,Gbit/s geostationary satellite path with an RTT of 600\,ms, without significantly
harming flows that shared the bottleneck. The authors also note that the pacing that
Careful Resume requires is challenging for UDP-based protocols such as QUIC. Unlike
TCP, which is paced inside the kernel, a QUIC stack runs in user space and has to pace
its packets itself, so the latency and jitter of system calls and operating system
scheduling make the gaps between packets less precise~\cite{kempf2025quicsteps}.

QUIC Delay Control (QUIC-DC)~\cite{mascolo2025quicdc}
estimates the one-way queueing delay of a connection to react to congestion early,
combines this with the TCP Westwood+~\cite{mascolo2001westwood} algorithm, and reduced losses and end-to-end delay compared with QUIC's CUBIC, BBRv2, NewReno and Westwood+, while preserving utilization.

Learning-based designs have been implemented in QUIC as well. PBQ~\cite{zhang2023pbq} combines BBR
with a reinforcement-learning agent: the agent sets the congestion window, while BBR
sets the pacing rate, and the combination achieved higher throughput and lower RTT
than QUIC with CUBIC or BBR. AERO~\cite{sidhu2026aero} is a learned congestion
controller designed to be integrated into any QUIC stack without extensive engineering
effort. For video on demand, it improved video quality (VMAF) by about 13\% and reduced
rebuffering by 30\% compared with existing approaches, and for low-latency live
streaming, it improved VMAF by about 12\% with 65\% fewer stalls.

Finally, Multipath QUIC lets one connection use several paths at once, for example
WiFi and LTE~\cite{deconinck2017mpquic}. XLINK~\cite{zheng2021xlink} applied it to short
videos, where it uses the video quality of experience perceived by the user
to control multipath scheduling.

\subsection{QUIC-specific improvements}\label{sec:rw-quic-specific}

Few works use features that TCP lacks, and most of them are not really congestion control
algorithms. They change what the controller knows when a connection starts, how losses
are detected and repaired, how often acknowledgments arrive, or how the capacity of a
connection is shared.

\paragraph{Path state in an extension frame.} The BDP\_FRAME~\cite{kuhn2021bdpframe} extension lets a QUIC
connection remember the characteristics of a path and use them when a session is
resumed. The server sends the client a frame with the RTT, the
congestion window and the client's IP address, and the client returns it when it
resumes the connection, so the server can start from the saved values. The address
lets the server detect whether the client's network has changed.

\paragraph{Client-provided hints.} Rüth et al.~\cite{ruth2019blitz} use the extensibility of QUIC to let a web client send the server an estimate of its bandwidth. The
server uses it to set the congestion window of a new connection, so the connection
skips slow start. The approach accepts some losses in exchange for a much shorter
FCT. For a typical DSL client, a 2\,MB video chunk loaded almost twice as fast, and in
the worst case, losses increased the transferred volume by 12\%.

\paragraph{Application information.} Wira~\cite{wu2024wira} sets the initial window and the initial rate
of a live-streaming connection from information about that
connection. It parses the first video frame and sizes the initial
window to fit it. Clients report quality-of-service metrics from earlier sessions in a
Transport Cookie, which the server uses to set the initial rate without storing state
for each client. Implemented on QUIC and deployed in commercial services, it reduced
the average time to complete the first frame by 10.6\%, and the 90th percentile by
16.7\%.

\paragraph{Loss repair.} QUIC-FEC~\cite{michel2019quicfec} adds a forward erasure correction frame to QUIC, so that lost packets can be rebuilt from redundant data without waiting for a
retransmission. To keep the loss signal for the congestion
controller, it distinguishes packets that were received from packets that were
recovered. FEC reduced the download time of smaller files on paths with high loss and
long delays, mainly by avoiding retransmission timeouts, but its overhead slowed down
long transfers and transfers on paths with little loss. FlEC~\cite{michel2023flec}
turns this into a framework inside QUIC, in which an application selects its
reliability mechanism, from pure retransmissions to different forms of FEC, and adapts
loss recovery and stream scheduling to its needs. For bulk transfers, file transfers
with restricted buffers and delay-constrained messages, it achieved lower latency than
QUIC's standard reliability mechanisms.

\paragraph{Acknowledgment frequency.} Custura et al.~\cite{custura2023reducing} showed that QUIC can send far
fewer acknowledgments without losing throughput, and proposed a policy for the
acknowledgment frequency extension: one
acknowledgment every ten packets, at least one every quarter of an RTT, and one every
two packets for the first 100 packets, to avoid slowing down slow
start.

% \paragraph{Handshake acknowledgments.} Because QUIC combines the transport and the
% cryptographic handshake, a server can acknowledge the client's first packet before it
% has its TLS certificate, an option called instant ACK. In common CDN
% deployments, the certificate has to be fetched from a backend store, so this choice
% determines the first RTT sample, from which the initial PTO is
% derived~\cite{mucke2024reacked}. M{\"u}cke et al.\ found that Cloudflare uses instant ACK
% for almost all of its domains. It shortens the initial PTO by three times the delay to
% the certificate store and speeds up recovery from losses during the handshake, but
% when that delay is high, it causes spurious retransmissions.

% \paragraph{Connection migration.} A QUIC connection can migrate to a new address, for
% example after a handover, without a new handshake. Chen et al.\ found that migration
% fails in up to 76\% of hard handovers, such as when a client moves from one Wi-Fi
% network to another, because QUIC starts a migration only passively, based on the
% timing of delivered data and on its loss detection. QUIC-HOA~\cite{chen2024quichoa} starts the migration
% actively instead, with a cross-layer design that is aware of the handover; it made
% migration succeed consistently and reduced the average migration time by up to
% 98\%.

% \paragraph{Real-time media over datagrams.} QUIC's unreliable DATAGRAM
% frames~\cite{rfc9221} allow real-time media to be carried inside a QUIC connection,
% which raises the question of how a media congestion controller interacts with the
% congestion control of QUIC. Engelbart and Ott implemented RTP over QUIC with GStreamer
% and quic-go and evaluated different combinations of congestion control algorithms and
% signaling~\cite{engelbart2021congestion}. In later work, they carried both the media
% and the data channels of WebRTC, which normally use separate protocols with separate
% congestion controllers, in a single QUIC connection. The media stream then saw lower
% latency and jitter, while the bandwidth was shared fairly with the data
% streams~\cite{engelbart2026quic}.

\paragraph{Stream-aware allocation.} The base QUIC protocol multiplexes many streams in
one connection, but offers no way to divide the bandwidth of the connection among them
by priority. PCC Priority~\cite{chen2023pccpriority} adds the priorities of the streams
to the utility function of PCC, so that the bandwidth of a connection is divided among
its streams in proportion to their priorities. In ns-3 simulations, it followed the priority
settings more closely than earlier schemes, allocating bandwidth at a ratio close to
that of the priorities, while keeping throughput high, latency low and its tolerance to
loss high. Unlike stream schedulers, which only decide which
stream sends next, it changes how the congestion control itself partitions the capacity
of the connection~\cite{rizvi2026cats}.

% CATS~\cite{rizvi2026cats} gives
% latency-critical messages priority over bulk data within a single TCP connection, using
% a scheduler and a credit-based shaper. Its author shows that the
% single in-order sequence space and the shared congestion window of TCP limit the tail
% latency that such a scheme can achieve, and proposes CATS over QUIC, with independent
% streams under one coupled congestion control, as the way around these limits.



